\section{LLM Embedding Models}
\medskip

Modern embedding models built from large language models solve a central weakness of earlier approaches such as Word2Vec. Word2Vec assigns a single vector $v_w$ to each word type $w$. This works surprisingly well for capturing general semantic similarity, but it cannot adapt to context. The word ``bank'' has one embedding whether we are talking about rivers or finance. The model has no mechanism for dynamically adjusting meaning.

Transformer-based language models change this completely. Instead of learning one vector per word type, they produce a representation for each \emph{token in context}. The embedding of a word is recomputed every time it appears. In effect, meaning is no longer stored statically in a dictionary. It is constructed dynamically from surrounding words.

You can think of Word2Vec as storing a fixed definition, while a transformer rewrites that definition each time the word appears, using the rest of the sentence as evidence.

\bigskip
\subsection{How Self-Attention Produces Context}
\medskip

The mechanism that enables contextual embeddings is self-attention. In each transformer layer, every input token embedding is linearly projected into three vectors: a query, a key, and a value. The queries are compared to keys using scaled dot products. These dot products determine attention weights, which specify how much each token should incorporate information from every other token.

The output for a token is a weighted combination of value vectors across the entire sequence. Because this process is repeated across multiple layers, information flows globally. Formally, in an encoder we can write
\[
h_t = f(x_1, \dots, x_T).
\]
The representation at position $t$ depends on the full input sequence.

This global interaction is what allows the embedding of ``wind'' in ``strong wind'' to differ from ``wind the clock.'' The surrounding words directly shape the vector representation.

\bigskip
\subsection{Encoders and Decoders}
\medskip

Transformer architectures differ primarily in how attention is structured.

In a decoder-only model, attention is causal. A token at position $t$ can only attend to positions $1, \dots, t$. Future tokens are masked:
\[
h_t = f(x_1, \dots, x_t).
\]
This structure enables next-word prediction and autoregressive generation. Models such as GPT are optimized for producing coherent text one token at a time.

In an encoder model, attention is non-causal. Every token can attend to every other token:
\[
h_t = f(x_1, \dots, x_T).
\]
Because each representation incorporates both left and right context, encoders are especially well suited for building embeddings that summarize entire inputs. For retrieval, semantic similarity, and classification, encoder-style models are typically preferred.

\bigskip

